% GENERATED FILE. Edit canonical lessons, then run npm run generate:books.
% canonical-source-hash: fnv1a64:3071c43e143a8294
% canonical-lessons: RU-C12-mama, RU-C12-papa, RU-C12-mat, RU-C12-otets

\chapter{Mother and Father}
\label{ch:russian-mother-and-father}
\hlchaptermodality{12}

\begin{chapteropening}
\textbf{By the end of this chapter:} \emph{I can name my mother and father in Russian, in both the affectionate and the formal register, and say which one is a real cousin of the English word and which is not.}

The chapter closes on \ru{отец}, the surprise of the pair: the reader produces \ru{мама}, \ru{папа}, \ru{мать} and \ru{отец}, states the natural-gender exception \ru{папа} tests, and names \ru{мать} as a secure PIE cousin of mother while \ru{отец} is not a cousin of father at all. It reaches back to Chapter 8's \ru{семья} to complete the family it gathers, and rescues four atoms nothing had revisited since Chapters 2, 7 and 8.
\end{chapteropening}

\section[máma]{\ru{мама} — \textquotedblleft{}mom,\textquotedblright{} a resemblance that isn't really a cousin}

\label{lesson:RU-C12-mama}



\begin{quote}
Say good night, say see-you-later, and now a new chapter — back to \textbf{\ru{семья}} (\emph{sem'yá}) to fill it in, one person at a time, starting with a word that looks like a cousin of English and, on close inspection, is not really one at all.
\end{quote}

\subsection*{You'll want to know first — \ru{мама}}
\begin{quote}
\textbf{\ru{мама}} — \emph{máma} — \textbf{mom, mum} (the everyday, affectionate word)
\end{quote}

Stress on the first syllable, \emph{MÁ-ma}. Ends in \textbf{-\ru{а}}, so the noun-gender ending rule holds: feminine, matching the person it names.

\begin{cousinweb}
Every etymology so far in this book has asked the same question: does this word share an \textbf{ancestor} with an English one? \emph{\ru{Мама}} is the first word where the honest answer requires a different question entirely.

\emph{\ru{Мама}} is not descended from Proto-Indo-European, and it is not borrowed from anywhere. Babies across the world, whatever language surrounds them, babble the same easy sounds first — an open \textbf{a}, made with a wide-open mouth, paired with the simplest consonants a mouth can make by pressing the lips together: \textbf{m}, \textbf{b}, \textbf{p}. Parents everywhere hear that babble and adopt it, so \emph{nearly every language on earth} ends up with something that sounds like \emph{mama} for \textquotedblleft{}mother\textquotedblright{} — Mandarin \textbf{māma}, French \textbf{maman}, Swahili \textbf{mama}, English \textbf{mama} — without any of them having borrowed the word or inherited it from a common source. It is not a family resemblance; it is a resemblance built from the shape of a human mouth, reinvented from scratch in every generation.

That is a genuinely different kind of \textquotedblleft{}no\textquotedblright{} from \emph{\ru{рот}}'s or \emph{\ru{до} \ru{скорого}}'s — those had no cousin \textbf{because} no one has found one. \emph{\ru{Мама}} has no cousin \textbf{because the question does not apply}: it was never inherited in the first place.
\end{cousinweb}

\subsection*{Your turn}
Say these aloud:

\begin{itemize}
\raggedright
  \item \textquotedblleft{}\ru{мама}\textquotedblright{} — feminine, first-syllable stress
  \item the babble sounds — \textquotedblleft{}open a, plus m, b, or p\textquotedblright{}
  \item the world tour — \textquotedblleft{}māma, maman, mama, mama\textquotedblright{}
  \item the distinction — \textquotedblleft{}not a family resemblance; the same invention, everywhere\textquotedblright{}
\end{itemize}

\begin{tcolorbox}[breakable,colback=teal!4,colframe=teal!35!black,title={Before you move on}]
Say \textquotedblleft{}mom.\textquotedblright{} (\textbf{\ru{Мама}}.) What gender, and why? (\textbf{Feminine} — ends in \textbf{-\ru{а}}.) Why does \emph{\ru{мама}} resemble the word for \textquotedblleft{}mother\textquotedblright{} in so many unrelated languages? (\textbf{Universal infant babbling} — the easiest sounds a baby's mouth makes, not shared ancestry.) Is this the same kind of \textquotedblleft{}no cousin\textquotedblright{} as \emph{\ru{рот}} or \emph{\ru{до} \ru{скорого}}? (\textbf{No} — those simply have not been traced to one; \emph{\ru{мама}} was never inherited at all.)
\end{tcolorbox}
\section[pápa]{\ru{папа} — \textquotedblleft{}dad,\textquotedblright{} the word that breaks the ending rule on purpose}

\label{lesson:RU-C12-papa}



\begin{quote}
\emph{\ru{Мама}}'s twin, built the same babbling way — but this one does something \emph{\ru{кофе}} only hinted at: it breaks the noun-gender ending rule outright, and tells you exactly when that is allowed to happen.
\end{quote}

\subsection*{You'll want to know first — \ru{папа}}
\begin{quote}
\textbf{\ru{папа}} — \emph{pápa} — \textbf{dad, papa} (the everyday, affectionate word)
\end{quote}

Stress on the first syllable, \emph{PÁ-pa} — the same rhythm as \emph{\ru{мама}}, the \textbf{\ru{м}} simply swapped for \textbf{\ru{п}}. Same babbling story, same easy vowel, one consonant apart.

\begin{grammarlens}[title={Grammar Lens: when biology beats the ending}]
Here is the rule you already know: \textbf{-\ru{а}/-\ru{я}} endings are feminine. \emph{\ru{Папа}} ends in \textbf{-\ru{а}} — and \emph{\ru{папа}} is \textbf{masculine}.

This is not another \emph{\ru{кофе}}, a fossil accident that happens to break the pattern. It is a real, live rule of its own, and a small closed set of words follows it: any noun naming a \textbf{male person} is masculine, whatever letter it ends in, because Russian lets biology outrank spelling for people. You will meet the rule's other members later — \textbf{\ru{дядя}} (\textquotedblleft{}uncle\textquotedblright{}) and \textbf{\ru{дедушка}} (\textquotedblleft{}grandpa\textquotedblright{}) both end in \textbf{-\ru{я}}/\textbf{-\ru{а}} and are both masculine for exactly this reason. \emph{\ru{Папа}} is your first sighting of it.
\end{grammarlens}

\subsection*{Your turn}
Say these aloud:

\begin{itemize}
\raggedright
  \item \textquotedblleft{}\ru{папа}\textquotedblright{} — masculine, despite the -\ru{а}
  \item the rule — \textquotedblleft{}a male person is always masculine, whatever the ending\textquotedblright{}
  \item the pair — \textquotedblleft{}\ru{мама}, feminine; \ru{папа}, masculine — same ending, opposite gender\textquotedblright{}
\end{itemize}

\begin{tcolorbox}[breakable,colback=teal!4,colframe=teal!35!black,title={Before you move on}]
Say \textquotedblleft{}dad.\textquotedblright{} (\textbf{\ru{Папа}}.) What gender is it, and does that match its \textbf{-\ru{а}} ending? (\textbf{Masculine} — it does \textbf{not} match; a male person overrides the ending.) Name the shared origin of \emph{\ru{мама}} and \emph{\ru{папа}}. (\textbf{Universal infant babbling} — the easiest mouth sounds, reinvented everywhere.) Name one other Russian word that will follow the same male-person rule later. (\textbf{\ru{Дядя}} or \textbf{\ru{дедушка}}.) And one last look further back: how do you say \textquotedblleft{}good night\textquotedblright{}? (\textbf{\ru{Спокойной} \ru{ночи}}.)
\end{tcolorbox}
\section[mat']{\ru{мать} — \textquotedblleft{}mother,\textquotedblright{} the word \ru{мама} only pretends to resemble}

\label{lesson:RU-C12-mat}



\begin{quote}
The formal word behind yesterday's affectionate one — and, unlike \emph{\ru{мама}}, a word with a real, unbroken, six-thousand-year paper trail straight to the English word you already own.
\end{quote}

\subsection*{You'll want to know first — \ru{мать}}
\begin{quote}
\textbf{\ru{мать}} — \emph{mat'} — \textbf{mother} (formal register: forms, introductions, \textquotedblleft{}my mother is a doctor\textquotedblright{})
\end{quote}

One syllable, softened by the \textbf{\ru{ь}} at the end — the same soft sign you already read in \emph{\ru{очень}} and \emph{\ru{семья}}. Feminine, matching \emph{\ru{мама}}, but reach for \emph{\ru{мать}} where English would reach for \emph{mother} rather than \emph{mom}: on a form, in an introduction, talking \emph{about} her rather than \emph{to} her.

\begin{cousinweb}
Where \emph{\ru{мама}} only \textbf{resembles} English \emph{mama} by babbling coincidence, \textbf{\ru{мать}} is the genuine article: Proto-Indo-European \textbf{*méh\textsubscript{2}tēr}, one of the oldest reconstructed words in the entire family, and one whose meaning has not shifted at all in perhaps six thousand years. English \textbf{mother}, Latin \textbf{māter} ($\to$ \emph{maternal}, \emph{matrix}), Greek \textbf{m\'{\={e}}tēr} ($\to$ \emph{metropolis}, literally \textquotedblleft{}mother city\textquotedblright{}), and Sanskrit \textbf{mātar} all continue the same word. Say \emph{\ru{мать}, mother, māter, m\'{\={e}}tēr, mātar} together and you are pronouncing one sentence's worth of unbroken history.

Linguists suspect *\emph{méh\textsubscript{2}tēr} itself grew out of the same infant babble as \emph{\ru{мама}} — the \emph{m}-sound a baby makes first — but so long ago, and so thoroughly grammaticalized into \textquotedblleft{}mother\textquotedblright{} across every daughter language, that it counts as real inheritance now, not live coincidence. \emph{\ru{Мама}} is that babble happening again, fresh, in front of you; \emph{\ru{мать}} is the fossil the same babble left five thousand years ago.
\end{cousinweb}

\subsection*{Your turn}
Say these aloud:

\begin{itemize}
\raggedright
  \item \textquotedblleft{}\ru{мать}\textquotedblright{} — one syllable, soft \ru{ть}
  \item \textquotedblleft{}\ru{мать}, mother, māter, m\'{\={e}}tēr, mātar\textquotedblright{} — one word, five languages
  \item the register split — \textquotedblleft{}\ru{мама} to her; \ru{мать} about her\textquotedblright{}
  \item the two kinds of \textquotedblleft{}m\textquotedblright{} — \textquotedblleft{}\ru{мама}, live babble; \ru{мать}, ancient fossil\textquotedblright{}
\end{itemize}

\begin{tcolorbox}[breakable,colback=teal!4,colframe=teal!35!black,title={Before you move on}]
Say \textquotedblleft{}mother,\textquotedblright{} formally. (\textbf{\ru{Мать}}.) Name three language cousins beyond English \emph{mother}. (\textbf{Latin \emph{māter}, Greek \emph{m\'{\={e}}tēr}, Sanskrit \emph{mātar}.}) Is \emph{\ru{мать}}'s connection to \emph{mother} the same kind of thing as \emph{\ru{мама}}'s? (\textbf{No} — \emph{\ru{мать}} is true inheritance; \emph{\ru{мама}} is fresh babble that merely resembles it.) When would you say \emph{\ru{мама}} instead of \emph{\ru{мать}}? (\textbf{To her, informally} — not on a form or in a formal introduction.)
\end{tcolorbox}
\section[atéts]{\ru{отец} — \textquotedblleft{}father,\textquotedblright{} an old word that is a stranger to its English namesake}

\label{lesson:RU-C12-otets}



\begin{quote}
\emph{\ru{Мать}} just gave you the cleanest cousin in this book. Its partner is the opposite kind of surprise: an equally ancient word — and one with \textbf{no} relation at all to the English word it translates.
\end{quote}

\subsection*{You'll want to know first — \ru{отец}}
\begin{quote}
\textbf{\ru{отец}} — \emph{atéts} — \textbf{father} (formal register, matching \emph{\ru{мать}})
\end{quote}

Stress on the last syllable, \emph{a-TÉTS}. Consonant-final, so — regular rule, no exception needed this time — masculine.

\begin{cousinweb}
Here is the twist \emph{\ru{мать}} did not prepare you for. \textbf{\ru{Отец}} is genuinely old — it continues Proto-Indo-European \textbf{*átta}, itself a nursery address-word for \textquotedblleft{}father,\textquotedblright{} parallel to \emph{\ru{мама}}'s babble but grammaticalized into the formal word millennia ago. Its cousins are real: \textbf{Gothic \emph{atta}}, \textbf{Latin \emph{atta}} (an affectionate, informal \textquotedblleft{}papa,\textquotedblright{} distinct from the formal \emph{pater}), \textbf{Hittite \emph{attaš}}, \textbf{Ancient Greek \emph{átta}.}

But \emph{none of those is English \textbf{father}}. English \emph{father} continues a \textbf{completely different} PIE root, \textbf{*ph\textsubscript{2}t\'{\={e}}r}, the formal word matched to \emph{*méh\textsubscript{2}tēr} \textquotedblleft{}mother\textquotedblright{} — the very cousin set \emph{\ru{мать}} just gave you. Slavic simply \textbf{lost} that word somewhere in its history and replaced it with the \emph{atta}-family word instead. So the two halves of this book's family pair are not symmetric at all: \textbf{\ru{мать}} kept the ancient formal word and still matches \emph{mother}; \textbf{\ru{отец}} kept an ancient word too, just a \emph{different} one, and the formal match with \emph{father} never survived into Slavic at all.

Introduce your own family the way you already learned to introduce yourself: \textbf{\ru{Меня} \ru{зовут}…} for you; \textbf{\ru{Его} \ru{зовут}…} (\textquotedblleft{}\textbf{his} name is…\textquotedblright{} — that same unnamed \textquotedblleft{}they,\textquotedblright{} and the same object-shaped pronoun \emph{\ru{его}} that \emph{\ru{вас}}/\emph{\ru{тебя}} already showed you a piece of) for your father, brother, or anyone else.
\end{cousinweb}

\subsection*{Your turn}
Say these aloud:

\begin{itemize}
\raggedright
  \item \textquotedblleft{}\ru{отец}\textquotedblright{} — \emph{a-TÉTS}, masculine
  \item the real cousins — \textquotedblleft{}atta, atta, attaš — Gothic, Latin, Hittite\textquotedblright{}
  \item the honest gap — \textquotedblleft{}NOT father — that word is \emph{ph\textsubscript{2}t\'{\={e}}r}, lost from Slavic entirely\textquotedblright{}
  \item the naming pattern — \textquotedblleft{}\ru{Меня} \ru{зовут}… / \ru{Его} \ru{зовут}…\textquotedblright{}
\end{itemize}

\begin{tcolorbox}[breakable,colback=teal!4,colframe=teal!35!black,title={Before you move on}]
Say \textquotedblleft{}father,\textquotedblright{} formally. (\textbf{\ru{Отец}}.) Is it related to English \emph{father}? (\textbf{No} — \emph{father} is \emph{*ph\textsubscript{2}t\'{\={e}}r}, a different root Slavic lost; \emph{\ru{отец}} continues \emph{*átta} instead.) Name three real cousins of \emph{\ru{отец}}. (\textbf{Gothic \emph{atta}, Latin \emph{atta}, Hittite \emph{attaš}.}) Now reach back to Chapters 2, 7 and 8: how would you say \textquotedblleft{}his name is…\textquotedblright{}, and what word gathers everyone this chapter just named? (\textbf{\ru{Его} \ru{зовут}…}; \textbf{\ru{семья}}, the word for \emph{\ru{друг}, \ru{подруга}, \ru{брат}, \ru{сестра}} — and now \emph{\ru{мама}, \ru{папа}, \ru{мать}, \ru{отец}} too.)
\end{tcolorbox}
